% Source: https://artofproblemsolving.com/community/c3h441901_help
% Stable identity: aops-topic-441901-q08
The base circumference $2\pi r=16\pi$ gives radius $r=8$ inches for both cones. The shorter cone has height $h$ with
\[
\dfrac{1}{3}\pi r^2h=\dfrac{1}{3}\pi\cdot64\,h=192\pi\qquad\Longrightarrow\qquad h=9 .
\]

\begin{asy}
usepackage("amsmath");
size(7cm);
defaultpen(fontsize(10pt));
real r=8;
void cone(pair c,real h,pen p){
  path front=c+(-r,0), back=c+(r,0);
  for(int i=1;i<=40;++i){ real t=pi+pi*i/40; front=front--(c+(r*cos(t),3*sin(t))); }
  for(int i=1;i<=40;++i){ real t=pi*i/40; back=back--(c+(r*cos(t),3*sin(t))); }
  draw(front,p);draw(back,p+dashed);
  draw(c+(-r,0)--c+(0,h)--c+(r,0),p);
  draw(c--c+(0,h),dashed);draw(c--c+(r,0),dashed);
  real s=0.9; draw(c+(s,0)--c+(s,s)--c+(0,s));
  label("$8$",c+(r/2,-1.3));
}
cone((0,0),30,linewidth(0.9));
cone((22,0),9,linewidth(0.9));
label("$30$",(0,15),W);label("$9$",(22,6),W);
label("original cone",(0,-5),S);label("shorter cone",(22,-5),S);
label("\mbox{}",(min(currentpicture,user=true)+max(currentpicture,user=true))/2,black); // keeps the text after the figure black when LaTeX typesets the labels inline
\end{asy}

The ratio of the heights is
\[
\dfrac{9}{30}=\boxed{\dfrac{3}{10}}.
\]
